\pdfbookmark[2]{On VI and VII.}{nexus-remarks-vi-vii}
\begin{center}\large On VI and VII.\end{center}

The way in which the author summed the series $\Sigma\left(\frac{-D}n\right)\frac1n$ can be seen from several particular examples found on separate sheets.

If $D\equiv3\pmod4$, the fundamental theorem in the theory of quadratic residues, together with the series
\[
\operatorname{cotang}u=\frac1u+\frac1{u-\pi}+\frac1{u+\pi}+\frac1{u-2\pi}+\frac1{u+2\pi}+\ldots
\]
gives
\[
\Sigma\left(\frac{-D}n\right)\frac1n=\Sigma\left(\frac nD\right)\frac1n=\frac\pi{2D}\Sigma\left(\frac\nu D\right)\operatorname{cotang}\frac{\nu\pi}{2D}
\]
where $\nu$ runs through all numbers relatively prime to $2D$ that are smaller than $D$. Put
\[
\surd-1=i,\quad\cos\frac{2\pi}D+i\sin\frac{2\pi}D=r
\]
and let $\mu$ denote all numbers relatively prime to $D$ that are no greater than $D$. The preceding sum can then readily be transformed into
\[
\Sigma\left(\frac{-D}n\right)\frac1n=\frac{\pi i}{4D}\left(\frac2D\right)\Sigma\left(\frac\mu D\right)\frac{r^\mu-1}{r^\mu+1}
\]
Now apply the formula valid for every root $\omega$ of the equation $\omega^D=1$,
\SourceLeaf{308}
\[
\frac{\omega-1}{\omega+1}=\Sigma(-1)^{\alpha-1}\omega^\alpha
\]
in which $\alpha$ must run through the numbers $1,2,3\ldots(D-1)$. Reversing the order of summation gives
\[
\Sigma\left(\frac{-D}n\right)\frac1n=\frac{\pi i}{4D}\left(\frac2D\right)\Sigma(-1)^{\alpha-1}\Sigma\left(\frac\mu D\right)r^{\alpha\mu}
\]
\Needspace{9\baselineskip}
The sum over $\mu$ can, as is known, be evaluated by means of the theorems proved in the paper \emph{Summatio quarundam serierum singularium}. Restricting attention to the case in which $D$ is divisible by no square, one finds in general
\begingroup\postdisplaypenalty=10000
\[
\Sigma\left(\frac\mu D\right)r^{\alpha\mu}=\left(\frac\alpha D\right)i^{\left(\frac{D-1}2\right)^2}\surd D
\]
\endgroup
where one must put $\left(\frac\alpha D\right)=0$ if $\alpha$ is not relatively prime to $D$. In the case $D\equiv3\pmod4$ one therefore obtains
\begingroup\postdisplaypenalty=10000\interlinepenalty=10000
\[
\Sigma\left(\frac{-D}n\right)\frac1n=\frac\pi{4\surd D}\left(\frac2D\right)\Sigma(-1)^\alpha\left(\frac\alpha D\right)=\frac\pi{2\surd D}\Sigma\left(\frac{\alpha^{\prime}}D\right)
\]
where $\alpha^{\prime}$ runs through all numbers relatively prime to $D$ that are smaller than $\frac12D$. Finally, since $\epsilon=2$, the number of classes becomes
\[
k=\Sigma\left(\frac{\alpha^{\prime}}D\right)
\]
\endgroup

If, on the other hand, $D\equiv1\pmod4$, the series
\[
\operatorname{cosec}u=\frac1u-\frac1{u-\pi}-\frac1{u+\pi}+\frac1{u-2\pi}+\frac1{u+2\pi}-\ldots
\]
gives in a similar way
\[
\begin{aligned}
\Sigma\left(\frac{-D}n\right)\frac1n&=\Sigma(-1)^{\frac{n-1}2}\left(\frac nD\right)\frac1n
=\frac\pi{2D}\Sigma(-1)^{\frac{\nu-1}2}\left(\frac\nu D\right)\operatorname{cosec}\frac{\nu\pi}{2D}\\
&=\frac\pi{2D}\Sigma\left(\frac\mu D\right)\frac{r^\mu}{r^{2\mu}+1}
\end{aligned}
\]
where the letters $\nu$ and $\mu$ have their earlier meanings. Excluding the evident case $D=1$, apply the formula valid for every root $\omega$ of the equation $\omega^D=1$ (except $\omega=1$),
\[
\frac\omega{\omega\omega+1}=1+\Sigma\omega^{4\alpha^{\prime\prime}}+\Sigma\omega^{D-4\alpha^{\prime\prime}}
\]
in which $\alpha^{\prime\prime}$ must run through the numbers $1,2,3\ldots\frac14(D-1)$. Again restricting $D$ to be divisible by no square, one obtains
\[
\Sigma\left(\frac{-D}n\right)\frac1n=\frac\pi{\surd D}\Sigma\left(\frac{\alpha^{\prime\prime}}D\right)
\]
and hence, since $\epsilon=2$,
\[
k=2\Sigma\left(\frac{\alpha^{\prime\prime}}D\right)
\]

\SourceLeaf{309}
The cases in which $D$ is even could be treated in a very similar way. —

\Needspace{19\baselineskip}
As for determining the number of classes for \emph{positive} determinants $D$, apart from the final formula given in the text, only a few geometrical figures survive. They depict hyperbolic sectors of finite dimensions, with inequalities beside them that confine the points whose coordinates are the variables of the quadratic forms to the interior of such a sector. These hyperbolic sectors take the place of the ellipses corresponding to quadratic forms of negative determinant. Determining their area again gives the average number of representations, provided that only representations whose variables satisfy those inequalities are admitted. On the other hand, the inequalities serve to isolate just one representation, excluding all others, from the infinitely many representations of a number $m$ that belong to one and the same root $n$ of the congruence $nn-D\equiv0\pmod{\frac m{\mu\mu}}$ and correspond to all the solutions of the equation $tt-Duu=1$ (compare Disqq. Arithm., art. 205). The total number of admitted representations of the number $m$ by the collection of all properly primitive forms that are not properly equivalent is then equal to the value of the function $f(m)$, in which only $-D$ is to be replaced by $D$. As before for negative determinants, considering its properties gives a second expression for the average number of representations; comparing this with the value previously derived by geometrical considerations then immediately determines the number of classes.
\BookRule{35mm}
